\section{Facts known only by computation}\label{sec:computed}

The statements of this section were verified by programs. They have no analytical proof here, and no statement of Sections~\ref{sec:complete}--\ref{sec:small} depends on them. For each item we name the program, which lies in the folder \texttt{note/checks/computed/} of the repository unless said otherwise, and say what it checks. We write $N_3(t,R)$ for the least order of a graph with chromatic number at least $t$ that has an $R$-compatible quantum $3$-coloring, as in~\cite{Z}.

\begin{itemize}
\item[(C1)] \emph{Criticality of $F_{22}$.} Deleting any vertex or any edge of $F_{22}$ leaves a graph with a proper $3$-coloring. Program: \texttt{checks/check\_small\_graph\_22.py} of~\cite{Z}, by exhaustive search over the $22+42$ deletions.

\item[(C2)] \emph{The $678$ edges of $F_*$ are necessary.} For each edge $\epsilon$ of $F_*$ there is a coloring of the $165$ measurements with four colors for which $\epsilon$ is the only edge of $F_\triangle$ with ends of equal color. Hence every subgraph of $F_\triangle$ with chromatic number five contains all edges of $F_*$; with Theorem~\ref{thm:star}, $F_*$ is the smallest subgraph of $F_\triangle$ with chromatic number five and the only one with $678$ edges, and it becomes $4$-colorable when any edge is deleted. Moreover, there are $678$ colorings of the $165$ measurements with four colors whose sets of edges of $F_{\rm phys}$ with ends of equal color are nonempty and pairwise disjoint. Hence every subgraph of $F_{\rm phys}$ with chromatic number five has at least $678$ edges. Program: \texttt{edges\_678/check\_678.py}, which rebuilds the three graphs and checks the colorings stored in \texttt{edge\_forced\_732.txt.gz} and \texttt{edge\_packing\_843.txt.gz}.

\item[(C3)] \emph{Smaller graphs for $R=2$ and $R=3$.} There is a graph with $207$ vertices, $1{,}071$ edges and chromatic number five that has a $2$-compatible quantum $3$-coloring in local dimension four, and a graph with $279$ vertices, $1{,}449$ edges and chromatic number five that has a $3$-compatible quantum $3$-coloring in local dimension four. Each becomes $4$-colorable when any vertex is deleted. Thus $N_3(5,2)\le207$ and $N_3(5,3)\le279$, compared with the $234$ and $288$ vertices of $F_2$ and $F_3$. Both graphs consist of six charts of $39$ words of the six bases, with lists of words and rules for the edges that differ from those of Section~\ref{sec:FR}. In the first graph three cells of each end grid are shared by the two charts of a wire. In the second, three of the nine wires join their two end grids directly and through three further vertices, and the other six through a layer of six vertices as in $F_3$. Each graph uses $165$ distinct measurements, not all of which belong to $F_{\rm phys}$. The lower bound for the chromatic number follows the parity argument of~\cite[Theorem~2.4]{Z}, with one step done by exhaustive search: every proper $4$-coloring of a modified chart reads the first coordinate, in the sense of~\cite[Section~2.4]{Z}, on an odd number of its three end grids. Programs: \texttt{radius\_two\_three/verify\_square\_radius\_two\_207.py}, which constructs the first graph with its projectors and checks the parity property, commutation for all $6{,}282$ pairs of vertices at distance at most two, a proper $5$-coloring and, with the option \texttt{--check-deletions square\_radius\_two\_207.json}, the stored $4$-colorings of the $207$ vertex-deleted subgraphs; and \texttt{verify\_square\_radius\_three\_replacements.py} with \texttt{verify\_replacement\_deletions.py} for the second graph. Independently of the parity argument, a search with a SAT solver confirms that neither graph has a proper $4$-coloring (\texttt{radius\_two\_three/sat\_crosscheck.py}, which is not part of the standard run because it takes about twenty minutes).

\item[(C4)] \emph{The graph $H_2$.} The graph $H_2$ of Theorem~\ref{thm:Hj} has $2{,}116{,}320{,}480$ edges. Let $\mathcal N$ be a set of first-level measurements that is independent in $F_{\rm phys}$, and let $\mathcal D$ be a family of nonzero projections on $\C^4$ any two of which have a nonzero product. The vertices $M$ of $H_2$ for which the sum of the projections $D_N(M)$, $N\in\mathcal N$, belongs to $\mathcal D$ form an independent set of $H_2$: for two such vertices some $D_N(M)D_{N'}(M')$ with $N,N'\in\mathcal N$ is nonzero, whereas adjacency would force $N\sim N'$ by the proof of Theorem~\ref{thm:Hj}(e). For the set $\mathcal N$ of $54$ measurements and the family $\mathcal D$ stored in \texttt{level\_two/independent\_set.json}, this independent set has $1{,}407{,}626$ vertices. Deleting it leaves a graph with $2{,}922{,}907$ vertices and $732{,}610{,}036$ edges. Its chromatic number is at least six, because a proper $5$-coloring would extend to a proper $6$-coloring of $H_2$, against Theorem~\ref{thm:Hj}; and it has a quantum $3$-coloring in local dimension $16$. Programs: \texttt{level\_two/family.py}, \texttt{count\_physical\_edges.py} and \texttt{count\_remaining\_edges.py}; they count by classes of vertices, and the second compares its method with a direct enumeration on small examples. They take the graph $F_{\rm phys}$ from the file \texttt{verify\_small\_graphs.py} of an earlier supplement of the project, which is included.
\end{itemize}

\paragraph{Facts with proofs that the programs replay.}
The program \texttt{note/checks/check\_note.py} recomputes the following facts, all of which are proved above: Lemma~\ref{lem:dictator} for $n=3$ and $L\le3$, by enumeration of all proper $4$-colorings of $K_3^{\times2}$ and $K_3^{\times3}$, and its parts~(a) and~(c) for $n=4$ and $L=2$, by enumeration of the $2{,}754{,}000$ proper $6$-colorings of $K_4^{\times2}$; the lower bound of Theorem~\ref{thm:complete} for the graph $S'_0(K_3)$, by exhaustive search; the identity~\eqref{eq:two-edges} on the $22$-vertex graph; the graphs $F_{22}^{(1)}$ and $F_{22}^{(2)}$ with their orders, sizes, chromatic number and coloring; the numbers $732$ and $678$, the table of Section~\ref{sec:165} and the contractibility of the $54$ complexes $\Lambda$ in twenty orders of deletion, and $\chi(F_*)\ge5$ by exhaustive search; the graphs $F_R$ for $R=2,3,4$ with their orders, the exact compatibility radius of the coloring and a proper $5$-coloring; and the count of the vertices of $H_2$, by enumeration over smaller sets of first-level measurements. One further test was attempted and did not finish: a search for a proof that $S_0(K_4)$ has no proper $6$-coloring, which would confirm the first case of Theorem~\ref{thm:complete} beyond~\cite{Z} without Lemma~\ref{lem:bg}, was stopped after nine minutes.
